\lampiransection{Dokumentasi Pra-Survei Observasi GrabFood}
\label{app:observasi-grabfood}

Lampiran ini menyajikan rekapitulasi data observasi awal pada halaman GrabFood Rumah Makan Nasi Gerilya. Pengamatan dilakukan pada tanggal 04 April 2026 untuk melihat informasi yang tampak bagi pelanggan, seperti tampilan toko, rating, jumlah penilaian, kisaran harga, menu terpopuler, dan promo yang tersedia. Sebagai pembanding awal, pengamatan juga mencatat dua usaha sejenis, yaitu RM Padang Paripurna dan RM Padang Dendeng Batokok. Informasi pembanding dan kecenderungan opini pelanggan digunakan pada Bab I untuk memperjelas konteks masalah penelitian.

Data pada lampiran ini digunakan sebagai gambaran awal kondisi Nasi Gerilya di GrabFood. Karena bersifat pra-survei, data tersebut belum dimaksudkan sebagai kesimpulan penelitian, melainkan sebagai pendukung latar belakang dan dasar penyusunan instrumen pengumpulan data.

\begingroup
\begin{xltabular}{\textwidth}{|Z{3.4cm}|Y|}
\caption{Identitas Observasi GrabFood}
\label{tab:identitas_observasi_grabfood} \\
        \toprule
        \textbf{Aspek}    & \textbf{Keterangan}                                                           \\ \hline
        \midrule
\endfirsthead
\continuedtabletitle{2}
        \toprule
        \textbf{Aspek}    & \textbf{Keterangan}                                                           \\ \hline
        \midrule
\endhead
\continuedtablefooter{2}
\endfoot
        Nama usaha        & Rumah Makan Nasi Gerilya                                                      \\ \hline
        Platform          & GrabFood                                                                      \\ \hline
        Tanggal observasi & 04 April 2026                                                                 \\ \hline
        Fokus observasi   & Tampilan toko, rating, penilaian, ulasan, harga, menu teratas, promo, dan opini konsumen \\ \hline
        \bottomrule
    
\end{xltabular}
\tablesource[\textwidth]{Observasi platform GrabFood pada 04 April 2026; diolah peneliti (2026).}
\endgroup

\begingroup
\begin{xltabular}{\textwidth}{|Z{3.4cm}|Y|}
\caption{Profil dan Data Utama Rumah Makan Nasi Gerilya}
\label{tab:data_gerilya} \\
        \toprule
        \textbf{Kategori Data} & \textbf{Fakta / Temuan}                                                                             \\ \hline
        \midrule
\endfirsthead
\continuedtabletitle{2}
        \toprule
        \textbf{Kategori Data} & \textbf{Fakta / Temuan}                                                                             \\ \hline
        \midrule
\endhead
\continuedtablefooter{2}
\endfoot
        Rating                 & 4,7/5,0                                                                                             \\ \hline
        Total penilaian        & 4.361 penilaian                                                                                     \\ \hline
        Harga terendah         & Rp11.000 (Nasi Putih)                                                                               \\ \hline
        Harga tertinggi        & Rp72.000 (Nasi Padang Kari Kambing Gerilya)                                                         \\ \hline
        Menu teratas           & Nasi Padang Rendang Sapi; Nasi Padang Ayam Goreng; Nasi Padang Ayam Bakar; Nasi Bakar Kikil Tunjang \\ \hline
        Harga menu teratas     & Rp56.000; Rp45.000; Rp45.000; Rp68.000                                                              \\ \hline
        Promo Grab             & Diskon Rp8.000 Jajan Hemat; Diskon 25\% maksimal Rp100.000; Diskon 15\% Group Order                 \\ \hline
        Promo bank             & Raya, BTN, BSI, Mega, Mega Syariah, SBI, OCBC, Danamon, Muamalat                                    \\ \hline
        \bottomrule
    
\end{xltabular}
\tablesource[\textwidth]{Observasi platform GrabFood pada 04 April 2026; diolah peneliti (2026).}
\endgroup

Daftar harga disusun dengan menelusuri menu yang tampil pada halaman GrabFood Rumah Makan Nasi Gerilya. Ringkasan harga setiap menu disajikan pada Tabel~\ref{tab:survei_harga_nasi_gerilya_grabfood}. 

\begingroup
\footnotesize
\begin{xltabular}{\textwidth}{|Z{2.4cm}|Z{3.2cm}|Y|}
\caption{Ringkasan Bukti Observasi GrabFood}
\label{tab:kode-bukti-prasurvei-grabfood} \\
        \toprule
        \textbf{Kode Data} & \textbf{Aspek} & \textbf{Keterangan} \\ \hline
        \midrule
\endfirsthead
\continuedtabletitle{3}
        \toprule
        \textbf{Kode Data} & \textbf{Aspek} & \textbf{Keterangan} \\ \hline
        \midrule
\endhead
\continuedtablefooter{3}
\endfoot
        K08-PB-001 & Rating dan penilaian & Menunjukkan nilai 4,7/5,0 dan jumlah penilaian awal Rumah Makan Nasi Gerilya pada GrabFood. \\ \hline
        PA-HRG-01 & Menu utama & Menunjukkan harga paket nasi Padang dan menu utama yang muncul pada halaman GrabFood. \\ \hline
        PA-HRG-02 & Lauk dan tambahan & Menunjukkan variasi lauk tambahan, pesanan tambahan, dan rentang harga menu pendamping. \\ \hline
        PA-HRG-03 & Sambal dan produk kemasan & Menunjukkan variasi sambal serta produk kemasan yang ditawarkan pada platform. \\ \hline
        PA-HRG-04 & Kuah dan minuman & Menunjukkan item tambahan berharga rendah serta menu minuman yang dapat menambah variasi transaksi. \\ \hline
        PA-HRG-05 & Menu premium & Menunjukkan keberadaan menu premium dengan harga tinggi pada halaman GrabFood. \\ \hline
        \bottomrule
    
\end{xltabular}
\tablesource[\textwidth]{Diolah peneliti dari observasi GrabFood tanggal 04 April 2026.}
\endgroup

\subsection*{Ringkasan Harga Menu}

Ringkasan harga menu berikut disajikan untuk mendokumentasikan data harga yang terlihat pada halaman GrabFood. 

% Data harga sudah terwakili oleh Tabel survei harga.
% \praSurveyEvidenceBlock{K08-PB-001}{Rating dan jumlah penilaian awal Rumah Makan Nasi Gerilya pada GrabFood.}{resource/03-prasurvey/grabfood-observation/cropped/K08-PB-001-rating-ulasan.jpeg}{0.60}
% \praSurveyEvidenceBlock{PA-HRG-01}{Tampilan menu utama dan harga paket nasi Padang pada halaman GrabFood.}{resource/03-prasurvey/grabfood-observation/cropped/PA-HRG-01-menu-utama.jpeg}{0.60}
% \praSurveyEvidenceBlock{PA-HRG-02}{Tampilan variasi lauk dan pesanan tambahan pada halaman GrabFood.}{resource/03-prasurvey/grabfood-observation/cropped/PA-HRG-02-pesanan-tambahan.jpeg}{0.60}
% \praSurveyEvidenceBlock{PA-HRG-03}{Tampilan variasi sambal dan produk kemasan pada halaman GrabFood.}{resource/03-prasurvey/grabfood-observation/cropped/PA-HRG-03-sambal-kemasan.jpeg}{0.55}
% \praSurveyEvidenceBlock{PA-HRG-04}{Tampilan menu kuah tambahan dan minuman pada halaman GrabFood.}{resource/03-prasurvey/grabfood-observation/cropped/PA-HRG-04-kuah-minuman.jpeg}{0.60}
% \praSurveyEvidenceBlock{PA-HRG-05}{Tampilan menu premium kepala ikan pada halaman GrabFood.}{resource/03-prasurvey/grabfood-observation/cropped/PA-HRG-05-menu-premium.jpeg}{0.48}

{\scriptsize
    
    

\refstepcounter{table}\label{tab:survei_harga_nasi_gerilya_grabfood}
    \edef\SurveiHargaNomor{\thetable}
    \newcommand{\BuktiHargaSatu}{PA-HRG-01}
    \newcommand{\BuktiHargaDua}{PA-HRG-02}
    \newcommand{\BuktiHargaTiga}{PA-HRG-03}
    \newcommand{\BuktiHargaEmpat}{PA-HRG-04}
    \newcommand{\BuktiHargaLima}{PA-HRG-05}
    \addcontentsline{lot}{table}{\protect\numberline{\thetable}Survei Harga Menu Rumah Makan Nasi Gerilya pada GrabFood}
    \begin{longtable}{|>{\raggedright\arraybackslash}p{\dimexpr\textwidth-4.8cm-6\tabcolsep-4\arrayrulewidth\relax}|>{\raggedleft\arraybackslash}p{2.4cm}|p{2.4cm}|}
        \multicolumn{3}{c}{\textbf{Tabel \SurveiHargaNomor}}                                        \\
        \multicolumn{3}{c}{\textbf{Survei Harga Menu Rumah Makan Nasi Gerilya pada GrabFood}}       \\[6pt]
        \toprule
        \textbf{Nama Menu}                                 & \textbf{Harga} & \textbf{Kode Data} \\ \hline
        \midrule
        \endfirsthead
        \continuedtabletitle{3}
        \toprule
        \textbf{Nama Menu}                                 & \textbf{Harga} & \textbf{Kode Data} \\ \hline
        \midrule
        \endhead
        \continuedtablefooter{3}
        \endfoot
        \bottomrule
        \endlastfoot
        

\rowcolor{black!8}\multicolumn{3}{|l|}{\textbf{A. GERILYA Nasi Padang}}                     \\ \hline
        Nasi Padang Dendeng Sapi Bakar GERILYA             & Rp56.000       & \BuktiHargaSatu       \\ \hline
        Nasi Padang Kikil Tunjang Gulai GERILYA            & Rp68.000       & \BuktiHargaSatu       \\ \hline
        Nasi Padang Ayam Goreng Bumbu GERILYA              & Rp45.000       & \BuktiHargaSatu       \\ \hline
        Nasi Padang Gulai Ayam Kalasan GERILYA             & Rp46.000       & \BuktiHargaSatu       \\ \hline
        Nasi Padang Ayam Bakar GERILYA                     & Rp45.000       & \BuktiHargaSatu       \\ \hline
        Nasi Padang Udang Goreng Panas GERILYA             & Rp58.000       & \BuktiHargaSatu       \\ \hline
        Nasi Padang Ikan Lele Goreng Panas GERILYA         & Rp39.000       & \BuktiHargaSatu       \\ \hline
        Nasi Padang Ikan Gembung Goreng GERILYA            & Rp50.000       & \BuktiHargaSatu       \\ \hline
        Nasi Padang Ikan Nila Goreng GERILYA               & Rp43.000       & \BuktiHargaSatu       \\ \hline
        Nasi Padang Ikan Nila Bakar Bumbu GERILYA          & Rp50.000       & \BuktiHargaSatu       \\ \hline
        Nasi Padang Ikan Kakap Goreng Panas GERILYA        & Rp48.000       & \BuktiHargaSatu       \\ \hline
        Nasi Padang Telur Gulai GERILYA                    & Rp32.000       & \BuktiHargaSatu       \\ \hline
        Nasi Padang Berkedel Kentang GERILYA               & Rp37.400       & \BuktiHargaSatu       \\ \hline
        Nasi Padang Sayur Polos GERILYA                    & Rp27.500       & \BuktiHargaSatu       \\ \hline
        Nasi Padang Gembung Kuring Bakar GERILYA           & Rp47.850       & \BuktiHargaSatu       \\ \hline
        Nasi Padang Kari Kambing GERILYA                   & Rp72.000       & \BuktiHargaSatu       \\ \hline
        Nasi Padang Ayam Pop Sauce Creamy GERILYA          & Rp46.000       & \BuktiHargaSatu       \\ \hline
        Nasi Padang Rendang Sapi GERILYA                   & Rp56.000       & \BuktiHargaSatu       \\ \hline
        Nasi Padang Ayam Rendang GERILYA                   & Rp46.000       & \BuktiHargaSatu       \\ \hline
        Nasi Padang Cumi Nagih GERILYA                     & Rp70.000       & \BuktiHargaSatu       \\ \hline
        Nasi Padang Ikan Marlin Goreng GERILYA             & Rp51.500       & \BuktiHargaDua        \\ \hline
        Nasi Padang Ikan Tongkol Tuna Goreng Panas GERILYA & Rp48.000       & \BuktiHargaDua        \\ \hline
        Nasi Padang Kakap Muncung Gulai Kuning GERILYA     & Rp51.150       & \BuktiHargaDua        \\ \hline
        Nasi Padang KAKAP MERAH Gulai Kuning GERILYA       & Rp69.300       & \BuktiHargaDua        \\ \hline
        Nasi Padang Telur Dadar GERILYA                    & Rp32.000       & \BuktiHargaDua        \\ \hline
        

\rowcolor{black!8}\multicolumn{3}{|l|}{\textbf{B. Pesanan Tambahan}}                        \\ \hline
        Jangek Siram                                       & Rp25.000       & \BuktiHargaDua        \\ \hline
        Keripik Kentang Balado                             & Rp27.500       & \BuktiHargaDua        \\ \hline
        Tambah Lauk Telur Gulai Merah                      & Rp17.500       & \BuktiHargaDua        \\ \hline
        Tambah Lauk Berkedel Kentang (size besar 85g)      & Rp17.500       & \BuktiHargaDua        \\ \hline
        Nasi Putih Tambah                                  & Rp11.000       & \BuktiHargaDua        \\ \hline
        Sate Rendang Jengkol                               & Rp7.700        & \BuktiHargaDua        \\ \hline
        Tambah Lauk Telur Dadar Sayur (1 butir lebih)      & Rp17.500       & \BuktiHargaDua        \\ \hline
        

\rowcolor{black!8}\multicolumn{3}{|l|}{\textbf{C. 1 Lauk, Ngak Cukup!}}                     \\ \hline
        Gulai cincang Gerilya                              & Rp55.000       & \BuktiHargaDua        \\ \hline
        Tambah Lauk Ayam Goreng Bumbu                      & Rp35.200       & \BuktiHargaDua        \\ \hline
        Tambah Lauk Gulai Ayam Kalasan                     & Rp36.000       & \BuktiHargaDua        \\ \hline
        Tambah Lauk Ayam Bakar Bumbu                       & Rp35.200       & \BuktiHargaDua        \\ \hline
        Tambah Lauk Gulai Kikil Tunjang                    & Rp63.000       & \BuktiHargaDua        \\ \hline
        Tambah Lauk Dendeng Sapi Bakar                     & Rp43.000       & \BuktiHargaDua        \\ \hline
        Tambah Lauk Udang Goreng Panas                     & Rp46.200       & \BuktiHargaDua        \\ \hline
        Tambah Lauk Ikan Lele Goreng Panas                 & Rp27.000       & \BuktiHargaDua        \\ \hline
        Tambah Lauk Ikan Nila Goreng Bumbu                 & Rp37.400       & \BuktiHargaTiga       \\ \hline
        Tambah Lauk Ikan Kakap Goreng Panas                & Rp38.000       & \BuktiHargaTiga       \\ \hline
        Tambah Lauk Kari Kambing                           & Rp67.500       & \BuktiHargaTiga       \\ \hline
        Tambah Lauk Ikan Gembung Bakar                     & Rp41.250       & \BuktiHargaTiga       \\ \hline
        Tambah Lauk Telur Ikan Sambal Merah GERILYA        & Rp40.150       & \BuktiHargaTiga       \\ \hline
        Tambah Lauk Ayam Pop Creamy GERILYA                & Rp38.500       & \BuktiHargaTiga       \\ \hline
        Tambah Lauk Ikan Marlin Goreng Bumbu               & Rp38.000       & \BuktiHargaTiga       \\ \hline
        Tambah Lauk Ayam Rendang (Kalasan)                 & Rp38.000       & \BuktiHargaTiga       \\ \hline
        Tambah Lauk Rendang Sapi                           & Rp46.000       & \BuktiHargaTiga       \\ \hline
        Tambah Lauk Cumi Nagih                             & Rp63.000       & \BuktiHargaTiga       \\ \hline
        Tambah Lauk Ikan Nila Bakar Bumbu                  & Rp41.000       & \BuktiHargaTiga       \\ \hline
        Tambah Lauk Ikan Gembung Kuring Goreng Panas       & Rp35.000       & \BuktiHargaTiga       \\ \hline
        Tambah Lauk Ikan Tongkol Tuna Goreng Bumbu         & Rp38.000       & \BuktiHargaTiga       \\ \hline
        Tambah Lauk KAKAP MERAH Gulai Kuning GERILYA       & Rp58.300       & \BuktiHargaTiga       \\ \hline
        Tambah Lauk Kakap Muncung Gulai Kuning GERILYA     & Rp40.150       & \BuktiHargaTiga       \\ \hline
        

\rowcolor{black!8}\multicolumn{3}{|l|}{\textbf{D. Sambal}}                                  \\ \hline
        GERILYA Sambal Hijau                               & Rp17.500       & \BuktiHargaTiga       \\ \hline
        GERILYA Sambal Merah                               & Rp17.500       & \BuktiHargaTiga       \\ \hline
        

\rowcolor{black!8}\multicolumn{3}{|l|}{\textbf{E. Sambal Gerilya Kemasan}}                  \\ \hline
        GERILYA Sambal Andaliman Getir (Kemasan)           & Rp38.500       & \BuktiHargaTiga       \\ \hline
        GERILYA Sambal Dencis Klotok (Kemasan)             & Rp38.500       & \BuktiHargaTiga       \\ \hline
        GERILYA Sambal Teri Medan Kali (Kemasan)           & Rp38.500       & \BuktiHargaEmpat      \\ \hline
        GERILYA Sambal Cumi Nagih (Kemasan)                & Rp38.500       & \BuktiHargaEmpat      \\ \hline
        GERILYA Sambal Tuna Asap Keumamah Aceh (Kemasan)   & Rp38.500       & \BuktiHargaEmpat      \\ \hline
        GERILYA Sambal Bawang Gurih (Kemasan)              & Rp38.500       & \BuktiHargaEmpat      \\ \hline
        

\rowcolor{black!8}\multicolumn{3}{|l|}{\textbf{F. Nambah Kuah}}                             \\ \hline
        Nambah KUAH CAMPUR                                 & Rp2.500        & \BuktiHargaEmpat      \\ \hline
        Nambah Kuah Gulai Kuning                           & Rp1.100        & \BuktiHargaEmpat      \\ \hline
        

\rowcolor{black!8}\multicolumn{3}{|l|}{\textbf{G. Minuman Segar}}                           \\ \hline
        Air Mineral (400ml)                                & Rp13.200       & \BuktiHargaEmpat      \\ \hline
        Es Teh Tawar                                       & Rp13.200       & \BuktiHargaEmpat      \\ \hline
        Es Teh Manis                                       & Rp15.950       & \BuktiHargaEmpat      \\ \hline
        Es Jeruk Limau                                     & Rp22.000       & \BuktiHargaEmpat      \\ \hline
        Es Lemon Tea                                       & Rp26.400       & \BuktiHargaEmpat      \\ \hline
        Kopi Hitam Dingin                                  & Rp23.100       & \BuktiHargaEmpat      \\ \hline
        Kopi Susu Dingin                                   & Rp26.400       & \BuktiHargaEmpat      \\ \hline
        Jus Jeruk Peras                                    & Rp22.000       & \BuktiHargaEmpat      \\ \hline
        Minuman Marqisa                                    & Rp20.000       & \BuktiHargaEmpat      \\ \hline
        

\rowcolor{black!8}\multicolumn{3}{|l|}{\textbf{H. Minuman Panas}}                           \\ \hline
        Teh Tawar Hangat                                   & Rp13.200       & \BuktiHargaEmpat      \\ \hline
        Teh Manis Hangat                                   & Rp15.950       & \BuktiHargaEmpat      \\ \hline
        Lemon Tea Panas                                    & Rp23.100       & \BuktiHargaEmpat      \\ \hline
        Kopi Hitam Panas                                   & Rp23.100       & \BuktiHargaEmpat      \\ \hline
        Kopi Susu Panas                                    & Rp26.400       & \BuktiHargaLima       \\ \hline
        

\rowcolor{black!8}\multicolumn{3}{|l|}{\textbf{I. GERILYA Gulai Merah Kepala Ikan}}         \\ \hline
        Gulai Merah Kepala Ikan GERILYA L                  & Rp275.000      & \BuktiHargaLima       \\ \hline
        Gulai Merah Kepala Ikan GERILYA M                  & Rp248.050      & \BuktiHargaLima       \\ \hline
    \end{longtable}
    \tablesource[\textwidth]{Observasi menu harga Rumah Makan Nasi Gerilya pada GrabFood, 04 April 2026; diolah peneliti (2026).}
    \addtocounter{table}{-1}
}

Data pada lampiran ini memperlihatkan kondisi awal Nasi Gerilya di GrabFood, terutama dari sisi tampilan toko, rentang harga, variasi menu, dan promosi yang tersedia. Uraian mengenai opini pelanggan, posisi terhadap kompetitor, dan arah analisis SWOT disampaikan pada bagian utama naskah serta akan diperdalam melalui pengumpulan data penelitian.

\subsection*{Tabel Harga Lengkap Nasi Gerilya}

Tabel berikut menyajikan daftar harga menu Nasi Gerilya secara lengkap, menggabungkan data pra-survei (04 April 2026) dan observasi utama (13 Juni 2026).

% Tabel harga lengkap Nasi Gerilya - dihasilkan dari data observasi

\begingroup
\scriptsize
\begin{xltabular}{\textwidth}{|C{0.7cm}|Y|R{1.8cm}|C{2.2cm}|}
\caption{Daftar Harga Menu Nasi Gerilya pada GrabFood}
\label{tab:daftar-harga-nasi-gerilya-grabfood} \\
\toprule
\textbf{No.} & \textbf{Nama Menu} & \textbf{Harga} & \textbf{Status} \\ \hline
\midrule
\endfirsthead

\continuedtabletitle{4}
\toprule
\textbf{No.} & \textbf{Nama Menu} & \textbf{Harga} & \textbf{Status} \\ \hline
\midrule
\endhead

\continuedtablefooter{4}
\endfoot

\bottomrule
\endlastfoot

% --- A. Menu Utama (Nasi Padang) ---


\rowcolor{black!8}\multicolumn{4}{|l|}{\textbf{A. Menu Utama (\textit{Nasi Padang})}} \\ \hline
1 & Nasi Padang Sayur Polos GERILYA & Rp27.500 & Tersedia \\ \hline
2 & Nasi Padang Telur Gulai GERILYA & Rp32.000 & Tersedia \\ \hline
3 & Nasi Padang Telur Dadar GERILYA & Rp32.000 & Tersedia \\ \hline
4 & Nasi Padang Berkedel Kentang GERILYA & Rp37.400 & Tersedia \\ \hline
5 & Nasi Padang Ikan Lele Goreng Panas GERILYA & Rp39.000 & Tersedia \\ \hline
6 & Nasi Padang Gulai Ayam Kalasan GERILYA & Rp41.500\textsuperscript{*} & Tersedia (Promo) \\ \hline
7 & Nasi Padang Ikan Nila Goreng GERILYA & Rp43.000 & Tersedia \\ \hline
8 & Nasi Padang Ayam Goreng Bumbu GERILYA & Rp45.000 & Tersedia (Terlaris) \\ \hline
9 & Nasi Padang Ayam Bakar GERILYA & Rp45.000 & Tersedia \\ \hline
10 & Nasi Padang Ayam Rendang GERILYA & Rp46.000 & Tersedia (Terlaris) \\ \hline
11 & Nasi Padang Ayam Pop Sauce Creamy GERILYA & Rp46.000 & Tidak Tersedia \\ \hline
12 & Nasi Padang Gembung Kuring Bakar GERILYA & Rp47.850 & Tidak Tersedia \\ \hline
13 & Nasi Padang Ikan Tongkol Tuna Goreng Panas GERILYA & Rp48.000 & Tersedia \\ \hline
14 & Nasi Padang Ikan Kakap Goreng Panas GERILYA & Rp48.000 & Tersedia \\ \hline
15 & Nasi Padang Ikan Gembung Goreng GERILYA & Rp50.000 & Tersedia \\ \hline
16 & Nasi Padang Bungkus Gulai Cincang GERILYA & Rp50.500\textsuperscript{*} & Tersedia (Promo) \\ \hline
17 & Nasi Padang Ikan Marlin Goreng GERILYA & Rp51.500 & Tersedia \\ \hline
18 & Nasi Padang Kakap Muncung Gulai Kuning GERILYA & Rp51.150 & Tersedia \\ \hline
19 & Nasi Padang Rendang Sapi GERILYA & Rp56.000 & Tersedia (Terlaris) \\ \hline
20 & Nasi Padang Dendeng Sapi Bakar GERILYA & Rp56.000 & Tersedia \\ \hline
21 & Nasi Padang Udang Goreng Panas GERILYA & Rp58.000 & Tersedia \\ \hline
22 & Nasi Padang Kikil Tunjang Gulai GERILYA & Rp68.000 & Tersedia \\ \hline
23 & Nasi Padang KAKAP MERAH Gulai Kuning GERILYA & Rp69.300 & Tersedia \\ \hline
24 & Nasi Padang Cumi Nagih GERILYA & Rp70.000 & Tersedia \\ \hline
25 & Nasi Padang Kari Kambing GERILYA & Rp72.000 & Tersedia \\ \hline
26 & Nasi Padang Ikan Nila Bakar Bumbu GERILYA & Rp50.000 & Tersedia \\ \hline

% --- B. Pesanan Tambahan ---


\rowcolor{black!8}\multicolumn{4}{|l|}{\textbf{B. Pesanan Tambahan}} \\ \hline
27 & Sate Rendang Jengkol & Rp7.700 & Tersedia \\ \hline
28 & Nasi Putih Tambah & Rp11.000 & Tersedia \\ \hline
29 & Tambah Lauk Telur Gulai Merah & Rp17.500 & Tersedia \\ \hline
30 & Tambah Lauk Berkedel Kentang & Rp17.500 & Tersedia \\ \hline
31 & Tambah Lauk Telur Dadar Sayur & Rp17.500 & Tersedia \\ \hline
32 & Jangek Siram & Rp25.000 & Tersedia \\ \hline
33 & Keripik Kentang Balado & Rp27.500 & Tersedia \\ \hline

% --- C. 1 Lauk, Ngak Cukup! (Lauk Tambahan) ---


\rowcolor{black!8}\multicolumn{4}{|l|}{\textbf{C. 1 Lauk, Ngak Cukup! (Lauk Tambahan)}} \\ \hline
34 & Gulai Cincang Gerilya & Rp55.000 & Tersedia \\ \hline
35 & Tambah Lauk Ikan Lele Goreng Panas & Rp27.000 & Tersedia \\ \hline
36 & Tambah Lauk Ikan Gembung Kuring Goreng Panas & Rp35.000 & Tersedia \\ \hline
37 & Tambah Lauk Ayam Goreng Bumbu & Rp35.200 & Tersedia \\ \hline
38 & Tambah Lauk Ayam Bakar Bumbu & Rp35.200 & Tersedia \\ \hline
39 & Tambah Lauk Ikan Nila Goreng Bumbu & Rp37.400 & Tersedia \\ \hline
40 & Tambah Lauk Ayam Pop Creamy GERILYA & Rp38.500 & Tersedia \\ \hline
41 & Tambah Lauk Ikan Kakap Goreng Panas & Rp38.000 & Tersedia \\ \hline
42 & Tambah Lauk Ikan Marlin Goreng Bumbu & Rp38.000 & Tersedia \\ \hline
43 & Tambah Lauk Ayam Rendang (Kalasan) & Rp38.000 & Tersedia \\ \hline
44 & Tambah Lauk Ikan Tongkol Tuna Goreng Bumbu & Rp38.000 & Tersedia \\ \hline
45 & Tambah Lauk Kakap Muncung Gulai Kuning GERILYA & Rp40.150 & Tersedia \\ \hline
46 & Tambah Lauk Telur Ikan Sambal Merah GERILYA & Rp40.150 & Tersedia \\ \hline
47 & Tambah Lauk Ikan Gembung Bakar & Rp41.250 & Tersedia \\ \hline
48 & Tambah Lauk Ikan Nila Bakar Bumbu & Rp41.000 & Tersedia \\ \hline
49 & Tambah Lauk Dendeng Sapi Bakar & Rp43.000 & Tersedia \\ \hline
50 & Tambah Lauk Gulai Ayam Kalasan & Rp36.000 & Tersedia \\ \hline
51 & Tambah Lauk Udang Goreng Panas & Rp46.200 & Tersedia \\ \hline
52 & Tambah Lauk Rendang Sapi & Rp46.000 & Tersedia \\ \hline
53 & Tambah Lauk KAKAP MERAH Gulai Kuning GERILYA & Rp58.300 & Tersedia \\ \hline
54 & Tambah Lauk Cumi Nagih & Rp63.000 & Tersedia \\ \hline
55 & Tambah Lauk Gulai Kikil Tunjang & Rp63.000 & Tersedia \\ \hline
56 & Tambah Lauk Kari Kambing & Rp67.500 & Tersedia \\ \hline

% --- D. Sambal ---


\rowcolor{black!8}\multicolumn{4}{|l|}{\textbf{D. Sambal}} \\ \hline
57 & GERILYA Sambal Hijau & Rp17.500 & Tersedia \\ \hline
58 & GERILYA Sambal Merah & Rp17.500 & Tersedia \\ \hline

% --- E. Sambal Gerilya Kemasan ---


\rowcolor{black!8}\multicolumn{4}{|l|}{\textbf{E. Sambal Gerilya Kemasan}} \\ \hline
59 & GERILYA Sambal Andaliman Getir (Kemasan) & Rp38.500 & Tersedia \\ \hline
60 & GERILYA Sambal Dencis Klotok (Kemasan) & Rp38.500 & Tersedia \\ \hline
61 & GERILYA Sambal Teri Medan Kali (Kemasan) & Rp38.500 & Tersedia \\ \hline
62 & GERILYA Sambal Cumi Nagih (Kemasan) & Rp38.500 & Tidak Tersedia \\ \hline
63 & GERILYA Sambal Tuna Asap Keumamah Aceh (Kemasan) & Rp38.500 & Tidak Tersedia \\ \hline
64 & GERILYA Sambal Bawang Gurih (Kemasan) & Rp38.500 & Tersedia \\ \hline

% --- F. Nambah Kuah ---


\rowcolor{black!8}\multicolumn{4}{|l|}{\textbf{F. Nambah Kuah}} \\ \hline
65 & Nambah Kuah Gulai Kuning & Rp1.100 & Tersedia \\ \hline
66 & Nambah KUAH CAMPUR & Rp2.500 & Tersedia \\ \hline

% --- G. Minuman Segar ---


\rowcolor{black!8}\multicolumn{4}{|l|}{\textbf{G. Minuman Segar}} \\ \hline
67 & Air Mineral (400ml) & Rp13.200 & Tersedia \\ \hline
68 & Es Teh Tawar & Rp13.200 & Tersedia \\ \hline
69 & Es Teh Manis & Rp15.950 & Tersedia \\ \hline
70 & Es Jeruk Limau & Rp22.000 & Tersedia \\ \hline
71 & Jus Jeruk Peras & Rp22.000 & Tersedia \\ \hline
72 & Minuman Marqisa & Rp20.000 & Tersedia \\ \hline
73 & Kopi Hitam Dingin & Rp23.100 & Tersedia \\ \hline
74 & Es Lemon Tea & Rp26.400 & Tersedia \\ \hline
75 & Kopi Susu Dingin & Rp26.400 & Tersedia \\ \hline

% --- H. Minuman Panas ---


\rowcolor{black!8}\multicolumn{4}{|l|}{\textbf{H. Minuman Panas}} \\ \hline
76 & Teh Tawar Hangat & Rp13.200 & Tersedia \\ \hline
77 & Teh Manis Hangat & Rp15.950 & Tersedia \\ \hline
78 & Lemon Tea Panas & Rp23.100 & Tersedia \\ \hline
79 & Kopi Hitam Panas & Rp23.100 & Tersedia \\ \hline
80 & Kopi Susu Panas & Rp26.400 & Tersedia \\ \hline

% --- I. Menu Premium Kepala Ikan ---


\rowcolor{black!8}\multicolumn{4}{|l|}{\textbf{I. Menu Premium Kepala Ikan}} \\ \hline
81 & Gulai Merah Kepala Ikan GERILYA M & Rp248.050 & Tidak Tersedia \\ \hline
82 & Gulai Merah Kepala Ikan GERILYA L & Rp275.000 & Tidak Tersedia \\ \hline

\end{xltabular}
\tablesource[\textwidth]{Observasi GrabFood Nasi Gerilya - Glugur Kota, 04 April 2026 dan 13 Juni 2026; diolah peneliti (2026).}

\noindent\footnotesize \textsuperscript{*}Harga setelah diskon promo. Harga normal: Gulai Ayam Kalasan Rp46.000; Gulai Cincang Rp55.000.
\endgroup


\subsection*{Tabel Harga Kompetitor}

Tabel berikut menyajikan daftar harga menu kelima kompetitor pembanding pada platform GrabFood, berdasarkan observasi tanggal 19 dan 20 Juni 2026.

\begingroup
\scriptsize

\begin{xltabular}{\textwidth}{|C{0.5cm}|Y|R{1.35cm}|C{0.5cm}|Y|R{1.35cm}|}
\caption{Daftar Menu dan Harga Kompetitor GrabFood}
\label{tab:harga-kompetitor-lengkap} \\
\toprule


\rowcolor{black!10}
\textbf{No.} & \textbf{Nama Menu} & \textbf{Harga} & \textbf{No.} & \textbf{Nama Menu} & \textbf{Harga} \\ \hline
\midrule
\endfirsthead

\continuedtabletitle{6}
\toprule


\rowcolor{black!10}
\textbf{No.} & \textbf{Nama Menu} & \textbf{Harga} & \textbf{No.} & \textbf{Nama Menu} & \textbf{Harga} \\ \hline
\midrule
\endhead

\continuedtablefooter{6}
\endfoot

\bottomrule
\endlastfoot



\rowcolor{black!12}\multicolumn{6}{|l|}{\textbf{Restoran Paripurna}} \\ \hline


\rowcolor{black!8}\multicolumn{6}{|l|}{\textbf{A. Menu Utama}} \\ \hline
1 & Ayam Gulai & Rp38.000 & 2 & Ayam Panggang & Rp38.000 \\ \hline
3 & Ayam Rendang & Rp38.000 & 4 & Ayam Goreng & Rp43.000 \\ \hline
5 & Ayam Bakar & Rp43.000 & 6 & Ayam Pop & Rp42.000 \\ \hline
7 & Rendang Daging & Rp44.000 & 8 & Dendeng Balado/Hijau & Rp44.000 \\ \hline
9 & Dendeng Lambok & Rp34.000 & 10 & Dendeng Batokok & Rp34.000 \\ \hline
11 & Ikan Kakap Goreng & Rp39.500 & 12 & Ikan Kakap Gulai & Rp39.500 \\ \hline
13 & Ikan Tenggiri & Rp36.000 & 14 & Ikan Tenggiri Rendang & Rp33.000 \\ \hline
15 & Nila Bakar & Rp37.200 & 16 & Gembung Bakar & Rp43.000 \\ \hline
17 & Gembung Asam Padeh & Rp43.000 & 18 & Tunjang & Rp50.500 \\ \hline
19 & Kari Kambing & Rp64.500 & 20 & Udang Toufu & Rp66.000 \\ \hline
21 & Sop Daging & Rp39.500 & & & \\ \hline


\rowcolor{black!8}\multicolumn{6}{|l|}{\textbf{B. Pelengkap dan Lauk Tambahan}} \\ \hline
22 & Nasi Bungkus Tambah & Rp7.000 & 23 & Kerupuk Udang & Rp10.000 \\ \hline
24 & Jangek & Rp11.500 & 25 & Keripik Kentang & Rp14.000 \\ \hline
26 & Puding & Rp14.000 & 27 & Nasi Putih & Rp15.000 \\ \hline
28 & Sambal Kentang Hati & Rp16.000 & 29 & Sambal Kecap & Rp18.000 \\ \hline
30 & Sambal Merah/Hijau & Rp18.000 & 31 & Telur Rendang & Rp18.000 \\ \hline
32 & Daun Ubi Tumbuk & Rp19.000 & 33 & Daun Ubi Rebus & Rp19.000 \\ \hline
34 & Sayur Lodeh & Rp19.000 & 35 & Buncis Tumis & Rp19.000 \\ \hline
36 & Perkedel & Rp24.300 & 37 & Telur Dadar & Rp24.300 \\ \hline
38 & Telur Gulai & Rp24.300 & 39 & Telur Sambal & Rp24.300 \\ \hline
40 & Jengkol Cabe Hijau & Rp28.500 & 41 & Jengkol Balado & Rp28.500 \\ \hline
42 & Jengkol Rendang & Rp28.500 & 43 & Nila Sambal Hijau & Rp28.500 \\ \hline
44 & Kakap Sambal & Rp30.000 & 45 & Gembung Acar & Rp31.500 \\ \hline
46 & Sambal Hati Puyuh & Rp31.500 & 47 & Buah Potong Biasa & Rp39.000 \\ \hline
48 & Sambal Cumi & Rp42.000 & 49 & Sambal Udang & Rp42.000 \\ \hline
50 & Sambal Kerang & Rp42.000 & 51 & Buah Potong Special & Rp47.000 \\ \hline
52 & Telur Gulai & Rp49.000 & 53 & Telur Sambal & Rp49.000 \\ \hline
54 & Perkedel & Rp53.000 & 55 & Sambal Hati Kentang & Rp53.000 \\ \hline
56 & Hati Kentang & Rp55.500 & 57 & Sambal Udang & Rp61.000 \\ \hline
58 & Jus Terong Belanda & Rp34.000 & & & \\ \hline


\rowcolor{black!8}\multicolumn{6}{|l|}{\textbf{C. Minuman}} \\ \hline
59 & Es Kosong & Rp5.000 & 60 & Teh Pahit Hangat & Rp8.000 \\ \hline
61 & Teh Pahit Dingin & Rp10.000 & 62 & Susu (Dingin/Hangat) & Rp10.000 \\ \hline
63 & Teh Botol Sosro & Rp10.000 & 64 & Air Mineral & Rp11.000 \\ \hline
65 & Teh Manis Panas & Rp11.000 & 66 & Teh Manis Dingin & Rp16.000 \\ \hline
67 & Teh Tarik (Dingin/Hangat) & Rp16.000 & 68 & Kopi Hitam & Rp16.000 \\ \hline
69 & Kopi Susu Biasa & Rp19.000 & 70 & Kopi Susu Paripurna & Rp20.000 \\ \hline
71 & Lemon Tea (Panas/Dingin) & Rp20.500 & 72 & Kelapa & Rp22.000 \\ \hline
73 & Jus Nanas & Rp24.300 & 74 & Timun Kerok & Rp24.300 \\ \hline
75 & Jus Tomat & Rp24.300 & 76 & Jus Pepaya & Rp24.300 \\ \hline
77 & Jus Belimbing & Rp24.300 & 78 & Jus Jeruk & Rp24.300 \\ \hline
79 & Jus Melon & Rp24.300 & 80 & Jus Semangka & Rp24.300 \\ \hline
81 & Jus Wortel & Rp24.500 & 82 & Jus Terong Belanda & Rp27.000 \\ \hline
83 & Jus Markisa & Rp27.000 & 84 & Jus Apel & Rp27.000 \\ \hline
85 & Jus Mangga & Rp27.000 & 86 & Jus Kuini & Rp27.000 \\ \hline
87 & Jus Guava & Rp27.000 & 88 & Jus Sirsak & Rp27.000 \\ \hline
89 & Jus Alpukat & Rp27.000 & 90 & Teh Serai Madu & Rp27.000 \\ \hline
91 & Jus Strawberry & Rp27.000 & 92 & Jus Mix & Rp34.000 \\ \hline
93 & Jus Worjer (Wortel + Jeruk) & Rp34.000 & 94 & Jus Wortom (Wortel + Tomat) & Rp34.000 \\ \hline
95 & Jus Awan (Jambu + Pepaya + Nanas) & Rp34.000 & 96 & Jus Bumi (Buah Naga + Mangga) & Rp34.000 \\ \hline
97 & Jus Matahari (Markisa + Terong Belanda + Jeruk) & Rp34.000 & 98 & Jus Sawi + Nanas & Rp34.000 \\ \hline
99 & Jus Buah Naga & Rp34.000 & 100 & Jus Durian & Rp34.000 \\ \hline
101 & Jus Kiwi & Rp34.000 & 102 & Jus Sunkist & Rp34.000 \\ \hline
103 & Jus Martabe (Markisa + Terong Belanda) & Rp34.000 & & & \\ \hline


\rowcolor{black!12}\multicolumn{6}{|l|}{\textbf{Rumah Makan Dendeng Batokok}} \\ \hline


\rowcolor{black!8}\multicolumn{6}{|l|}{\textbf{A. Menu Utama}} \\ \hline
104 & Lele Goreng & Rp12.000 & 105 & Lele Bakar & Rp12.000 \\ \hline
106 & Kikil & Rp13.000 & 107 & Ikan Tongkol Asam Padeh & Rp13.000 \\ \hline
108 & Ikan Tongkol Goreng & Rp13.000 & 109 & Ikan Dencis Goreng & Rp13.000 \\ \hline
110 & Ayam Saos & Rp15.000 & 111 & Ayam Kecap & Rp15.000 \\ \hline
112 & Ayam Bakar & Rp15.000 & 113 & Ayam Kemangi & Rp15.000 \\ \hline
114 & Ayam Rendang & Rp15.000 & 115 & Ayam Goreng & Rp15.000 \\ \hline
116 & Kepala Nila & Rp16.000 & 117 & Ikan Nila Goreng & Rp16.000 \\ \hline
118 & Ikan Nila Bakar & Rp16.000 & 119 & Ikan Nila Ijo & Rp16.000 \\ \hline
120 & Nasi Lele Goreng & Rp17.000 & 121 & Nasi Lele Bakar & Rp17.000 \\ \hline
122 & Gulai & Rp18.000 & 123 & Ikan Gembung Bakar & Rp18.000 \\ \hline
124 & Ikan Kakap Goreng & Rp18.000 & 125 & Nasi Kikil & Rp18.000 \\ \hline
126 & Nasi Tongkol Asam Padeh & Rp18.000 & 127 & Nasi Tongkol Goreng & Rp18.000 \\ \hline
128 & Nasi Dencis Goreng & Rp18.000 & 129 & Nasi Ayam Goreng Tepung & Rp20.000 \\ \hline
130 & Nasi Ayam Saos & Rp20.000 & 131 & Nasi Ayam Kecap & Rp20.000 \\ \hline
132 & Nasi Ayam Kemangi & Rp20.000 & 133 & Nasi Ayam Cabe Ijo & Rp20.000 \\ \hline
134 & Nasi Ayam Bakar & Rp20.000 & 135 & Nasi Ayam Goreng & Rp20.000 \\ \hline
136 & Nasi Gulai Kepala Nila & Rp21.000 & 137 & Nasi Nila Goreng & Rp21.000 \\ \hline
138 & Nasi Nila Bakar & Rp21.000 & 139 & Nasi Nila Ijo & Rp21.000 \\ \hline
140 & Babat & Rp21.500 & 141 & Nasi Usus & Rp22.500 \\ \hline
142 & Rendang Daging & Rp23.000 & 143 & Nasi Gembung Bakar & Rp23.000 \\ \hline
144 & Nasi Kakap Gulai & Rp23.000 & 145 & Cincang & Rp25.000 \\ \hline
146 & Nasi Rendang & Rp29.000 & 147 & Nasi Dendeng Basah & Rp29.000 \\ \hline
148 & Gulai Tunjang & Rp30.000 & 149 & Nasi Cincang & Rp32.000 \\ \hline
150 & Nasi Gulai Tunjang & Rp36.000 & & & \\ \hline


\rowcolor{black!8}\multicolumn{6}{|l|}{\textbf{B. Pelengkap dan Lauk Tambahan}} \\ \hline
151 & Perkedel & Rp7.000 & 152 & Telur Gulai & Rp7.000 \\ \hline
153 & Sambal & Rp7.000 & 154 & Puding & Rp9.000 \\ \hline
155 & Kerupuk & Rp10.000 & 156 & Kerupuk Jangek & Rp10.000 \\ \hline
157 & Telur Dadar & Rp10.000 & 158 & Terong Sambal & Rp11.000 \\ \hline
159 & Nasi Sayur & Rp12.000 & 160 & Sambal Teri Kacang & Rp13.000 \\ \hline
161 & Kentang Hati & Rp13.000 & 162 & Kentang Puyuh & Rp13.000 \\ \hline
163 & Jengkol Sambal & Rp13.000 & 164 & Buah Potong & Rp14.000 \\ \hline
165 & Ayam Sambal Merah & Rp15.000 & 166 & Nasi Tahu+Tempe & Rp15.000 \\ \hline
167 & Nasi Perkedel & Rp15.000 & 168 & Nasi Telur Sambal & Rp15.000 \\ \hline
169 & Nasi Telur Gulai & Rp15.000 & 170 & Nasi Telur Dadar & Rp17.000 \\ \hline
171 & Ikan Gembung Acar Rica & Rp18.000 & 172 & Nasi Kentang Hati & Rp18.000 \\ \hline
173 & Nasi Kentang Puyuh & Rp18.000 & 174 & Nasi Jengkol Sambal & Rp18.000 \\ \hline
175 & Nasi Ayam Sambal Merah & Rp20.000 & 176 & Udang Sambal & Rp21.000 \\ \hline
177 & Dendeng Sambal & Rp23.000 & 178 & Nasi Gembung Acar Rica & Rp23.000 \\ \hline
179 & Nasi Kakap Goreng Sambal & Rp23.000 & 180 & Udang Sambal & Rp26.500 \\ \hline
181 & Kentang Hati & Rp18.000 & & & \\ \hline


\rowcolor{black!8}\multicolumn{6}{|l|}{\textbf{C. Minuman}} \\ \hline
182 & Teh Pahit Dingin & Rp5.000 & 183 & Air Mineral Prima & Rp8.500 \\ \hline
184 & Teh Manis Dingin & Rp9.000 & 185 & Fruit Tea & Rp10.000 \\ \hline
186 & Teh Botol Sosro & Rp10.000 & 187 & Jus Nenas & Rp16.500 \\ \hline
188 & Jus Wortel & Rp16.500 & 189 & Jus Timun & Rp16.500 \\ \hline
190 & Jus Kuini & Rp16.500 & 191 & Jus Sirsak & Rp16.500 \\ \hline
192 & Jus Terong Belanda & Rp16.500 & 193 & Jus Buah Naga & Rp16.500 \\ \hline
194 & Jus Jeruk & Rp16.500 & 195 & Jus Mangga & Rp16.500 \\ \hline
196 & Jus Alpukat & Rp16.500 & 197 & Jus Apel Hijau & Rp16.500 \\ \hline


\rowcolor{black!12}\multicolumn{6}{|l|}{\textbf{Rumah Makan Istana Krakatau}} \\ \hline


\rowcolor{black!8}\multicolumn{6}{|l|}{\textbf{A. Menu Utama}} \\ \hline
198 & Ayam Bakar & Rp47.000 & 199 & Dendeng Sapi & Rp48.000 \\ \hline
200 & Rendang Sapi & Rp48.000 & 201 & Kakap Goreng & Rp48.000 \\ \hline
202 & Kakap Gulai & Rp50.000 & 203 & Kakap Kukus & Rp50.000 \\ \hline
204 & Nasi Dendeng Sapi & Rp47.000 & 205 & Nasi Rendang Sapi & Rp63.000 \\ \hline
206 & Nasi Ayam Bakar & Rp63.000 & 207 & Nasi Kotak Ayam Gulai & Rp65.000 \\ \hline
208 & Nasi Kotak Ayam Rendang & Rp65.000 & 209 & Nasi Kotak Ayam Goreng & Rp65.000 \\ \hline
210 & Nasi Kotak Rendang Sapi & Rp67.000 & 211 & Nasi Kotak Ikan & Rp67.000 \\ \hline
212 & Nasi Kotak Dendeng Sapi & Rp67.000 & 213 & Nasi Kotak Ikan Gulai & Rp72.000 \\ \hline
214 & Belut Cabe Ijo & Rp98.000 & 215 & Cumi Tauco & Rp105.000 \\ \hline
216 & Cumi Goreng Tepung & Rp108.000 & 217 & Ayam Goreng Per Ekor & Rp133.000 \\ \hline
218 & Bawal Kukus & Rp155.000 & 219 & Ayam Saus Istana & Rp198.000 \\ \hline
220 & Udang Gala Goreng & Rp213.000 & 221 & Udang Gala Lada Hitam & Rp223.000 \\ \hline
222 & Udang Gala Saus Padang & Rp223.000 & 223 & Udang Tauco & Rp240.000 \\ \hline
224 & Udang Kelong Goreng & Rp240.000 & 225 & Udang Kelong Mentega & Rp240.000 \\ \hline
226 & Udang Kelong Cabe Ijo & Rp240.000 & 227 & Udang Kelong Sambal & Rp240.000 \\ \hline
228 & Bawal Saus Padang & Rp250.000 & 229 & Kepala Ikan Gulai & Rp298.000 \\ \hline
230 & Teri Joss & Rp42.000 & & & \\ \hline


\rowcolor{black!8}\multicolumn{6}{|l|}{\textbf{B. Pelengkap dan Lauk Tambahan}} \\ \hline
231 & Perkedel Kentang & Rp19.000 & 232 & Telur Mata Sapi & Rp19.000 \\ \hline
233 & Telur Bulat & Rp19.000 & 234 & Opak & Rp20.000 \\ \hline
235 & Telur Dadar & Rp21.000 & 236 & Nasi Putih & Rp22.500 \\ \hline
237 & Tempe Goreng Per Porsi & Rp25.000 & 238 & Tahu Goreng Per Porsi & Rp28.000 \\ \hline
239 & Bunga Pepaya & Rp32.000 & 240 & Genjer & Rp32.000 \\ \hline
241 & Urap & Rp32.000 & 242 & Acar Timun & Rp32.000 \\ \hline
243 & Daun Ubi Tumbuk & Rp32.000 & 244 & Buncis & Rp32.000 \\ \hline
245 & Terong Belanda & Rp35.000 & 246 & Sambal Tempe & Rp37.000 \\ \hline
247 & Pete Sambal & Rp42.000 & 248 & Terong Sambal & Rp42.000 \\ \hline
249 & Rendang Jengkol & Rp45.000 & 250 & Perkedel Jagung & Rp46.000 \\ \hline
251 & Kakap Sambal & Rp48.000 & 252 & Nasi Telur Dadar & Rp51.000 \\ \hline
253 & Nasi Ikan Cabe Ijo & Rp65.000 & 254 & Nasi Kotak Ayam Sambal & Rp65.000 \\ \hline
255 & Ayam Cabe Ijo Per Porsi & Rp93.000 & 256 & Dendeng Batokok Per Porsi & Rp95.000 \\ \hline
257 & Udang Kelong Cabe Ijo & Rp240.000 & 258 & Udang Kelong Sambal & Rp240.000 \\ \hline


\rowcolor{black!8}\multicolumn{6}{|l|}{\textbf{C. Minuman}} \\ \hline
259 & Teh Tong & Rp9.000 & 260 & Teh Pahit Dingin & Rp10.000 \\ \hline
261 & Air Mineral & Rp11.000 & 262 & Teh Manis Panas & Rp15.000 \\ \hline
263 & Teh Manis Dingin & Rp15.000 & 264 & Cincau & Rp30.000 \\ \hline
265 & Lemon Tea & Rp30.000 & 266 & Jus Tomat & Rp35.000 \\ \hline
267 & Jeruk & Rp35.000 & 268 & Kelapa Thai & Rp35.000 \\ \hline
269 & Jus Timun & Rp35.000 & 270 & Jus Sirsak & Rp35.000 \\ \hline
271 & Jus Murni Tanpa Es & Rp38.000 & & & \\ \hline


\rowcolor{black!12}\multicolumn{6}{|l|}{\textbf{RM. Pondok Gurih}} \\ \hline


\rowcolor{black!8}\multicolumn{6}{|l|}{\textbf{A. Menu Utama}} \\ \hline
272 & Nasi + Ayam Goreng & Rp50.000 & 273 & Nasi + Ayam Rendang & Rp50.000 \\ \hline
274 & Nasi + Ayam Gulai & Rp50.000 & 275 & Ayam Goreng & Rp35.750 \\ \hline
276 & Ayam Gulai & Rp35.750 & 277 & Ayam Rendang & Rp35.750 \\ \hline
278 & Ayam Bakar & Rp35.750 & 279 & Ikan Goreng & Rp35.750 \\ \hline
280 & Ikan Gulai & Rp35.750 & 281 & Ikan Gembung Bakar & Rp35.750 \\ \hline
282 & Ikan Asam Pedas & Rp35.750 & 283 & Rendang Sapi & Rp43.250 \\ \hline
284 & Dendeng & Rp43.250 & 285 & Kikil & Rp43.250 \\ \hline
286 & Teri Cabai Hijau & Rp47.500 & 287 & Ikan Bawal Goreng & Rp50.000 \\ \hline
288 & Gembung Bakar & Rp50.000 & 289 & Kari Kambing & Rp64.500 \\ \hline
290 & Kepala Ikan Kecil & Rp180.000 & 291 & Kepala Ikan Sedang & Rp213.500 \\ \hline
292 & Kepala Ikan Besar & Rp246.750 & & & \\ \hline


\rowcolor{black!8}\multicolumn{6}{|l|}{\textbf{B. Pelengkap dan Lauk Tambahan}} \\ \hline
293 & Kerupuk Emping & Rp15.000 & 294 & Kerupuk Jangek & Rp15.000 \\ \hline
295 & Telur Bulat/Sambal & Rp18.500 & 296 & Telur Mata Sapi & Rp18.500 \\ \hline
297 & Perkedel & Rp19.750 & 298 & Telur Dadar & Rp19.750 \\ \hline
299 & Sayur Daun Ubi Rebus & Rp19.750 & 300 & Sambal Hati Ampela & Rp22.000 \\ \hline
301 & Ikan Sambal & Rp35.750 & 302 & Ikan Gembung Acar & Rp35.750 \\ \hline
303 & Sayur Gori & Rp35.750 & 304 & Sayur Tauco & Rp35.750 \\ \hline
305 & Sayur Daun Ubi Gulai & Rp35.750 & 306 & Sayur Lodeh & Rp35.750 \\ \hline
307 & Sayur Daun Ubi Tumbuk & Rp35.750 & 308 & Nasi Sayur & Rp37.500 \\ \hline
309 & Telur Bulat & Rp42.250 & 310 & Telur Mata Sapi & Rp42.250 \\ \hline
311 & Perkedel & Rp42.250 & 312 & Telur Dadar & Rp42.250 \\ \hline
313 & Sambal Goreng Kentang & Rp43.500 & 314 & Jengkol & Rp47.500 \\ \hline
315 & Ikan Bawal Sambal & Rp50.000 & 316 & Ikan Sambal & Rp50.000 \\ \hline
317 & Gembung Acar & Rp50.000 & 318 & Sambal Udang Pete & Rp68.500 \\ \hline


\rowcolor{black!8}\multicolumn{6}{|l|}{\textbf{C. Minuman}} \\ \hline
319 & Teh Tong & Rp6.750 & 320 & Teh Pahit Dingin & Rp8.000 \\ \hline
321 & Le Minerale 600 Ml & Rp9.500 & 322 & Teh Botol & Rp10.750 \\ \hline
323 & Fruit Tea & Rp10.750 & 324 & Le Minerale 1500 Ml & Rp12.000 \\ \hline
325 & Teh Manis Dingin & Rp14.750 & 326 & Teh Manis Hangat & Rp14.750 \\ \hline
327 & Badak & Rp16.250 & 328 & Es Sarang Burung & Rp18.500 \\ \hline
329 & Teh Bunga (Liang Teh) & Rp20.500 & 330 & Kelapa Muda & Rp31.500 \\ \hline
331 & Jus Guava & Rp31.750 & 332 & Jus Timun & Rp31.750 \\ \hline
333 & Jus Tomat & Rp31.750 & 334 & Jus Sirsak & Rp31.750 \\ \hline
335 & Jus Jeruk & Rp31.750 & 336 & Jus Sunkist & Rp31.750 \\ \hline
337 & Jus Belimbing & Rp31.750 & 338 & Jus Alpukat & Rp34.250 \\ \hline
339 & Jus Terong Belanda & Rp37.000 & 340 & Jus Jeruk Murni & Rp37.000 \\ \hline
341 & Jus Mangga & Rp37.000 & 342 & Jus Martabe & Rp37.000 \\ \hline
343 & Jus Kuini & Rp37.000 & & & \\ \hline


\rowcolor{black!12}\multicolumn{6}{|l|}{\textbf{Restoran Garuda}} \\ \hline


\rowcolor{black!8}\multicolumn{6}{|l|}{\textbf{A. Menu Utama}} \\ \hline
344 & Nasi Bks + Ayam Goreng & Rp49.206 & 345 & Nasi Bks + Ayam Rendang & Rp49.206 \\ \hline
346 & Nasi Bks + Ayam Panggang & Rp49.206 & 347 & Nasi Bks + Ayam Sambal Ijo & Rp49.206 \\ \hline
348 & Nasi Bks + Ayam Saos & Rp49.206 & 349 & Nasi Bks + Ayam Semur & Rp49.206 \\ \hline
350 & Nasi Bks + Ayam + Telur & Rp68.254 & 351 & Ikan Lele & Rp36.508 \\ \hline
352 & Usus & Rp38.095 & 353 & Babat & Rp38.095 \\ \hline
354 & Ayam Semur & Rp39.683 & 355 & Ayam Panggang & Rp39.683 \\ \hline
356 & Rendang Limpa & Rp41.270 & 357 & Rendang Hati & Rp41.270 \\ \hline
358 & Ikan Gembung Panggang & Rp41.270 & 359 & Paru & Rp41.270 \\ \hline
360 & Gado-Gado & Rp41.270 & 361 & Ikan Gulai Patin & Rp42.251 \\ \hline
362 & Ikan Goreng & Rp42.857 & 363 & Ikan Gulai & Rp42.857 \\ \hline
364 & Dendeng Kering & Rp46.032 & 365 & Ayam Rendang & Rp49.206 \\ \hline
366 & Ayam Goreng & Rp49.206 & 367 & Ayam Saos & Rp49.206 \\ \hline
368 & Ayam Semur & Rp49.206 & 369 & Hati & Rp49.206 \\ \hline
370 & Gembung Panggang & Rp50.794 & 371 & Rendang & Rp52.381 \\ \hline
372 & Ikan Kakap Gulai & Rp52.381 & 373 & Dendeng & Rp53.968 \\ \hline
374 & Ayam Gulai & Rp55.556 & 375 & Gulai Otak & Rp57.143 \\ \hline
376 & Kikil & Rp58.730 & 377 & Rendang Daging & Rp58.730 \\ \hline
378 & Kakap Goreng & Rp58.730 & 379 & Cumi & Rp66.667 \\ \hline
380 & Sate Padang & Rp66.667 & 381 & Sate Ayam (Bumbu Kacang) & Rp66.667 \\ \hline
382 & Kerang & Rp73.016 & 383 & Ikan Teri Bilis & Rp74.603 \\ \hline
384 & Martabak Mesir & Rp74.603 & 385 & Udang Goreng Nestum & Rp76.051 \\ \hline
386 & Teri Bilis & Rp80.952 & 387 & Kari Kambing & Rp85.714 \\ \hline
388 & Sop Daging & Rp88.889 & 389 & Ikan Bawal Steam & Rp113.751 \\ \hline


\rowcolor{black!8}\multicolumn{6}{|l|}{\textbf{B. Pelengkap dan Lauk Tambahan}} \\ \hline
390 & Gulai Tahu & Rp14.300 & 391 & Gulai Tempe & Rp14.300 \\ \hline
392 & Nasi Putih Bungkus & Rp19.048 & 393 & Hanya Lauk Saja & Rp22.222 \\ \hline
394 & Telur Dadar & Rp25.397 & 395 & Perkedel Jagung & Rp25.397 \\ \hline
396 & Perkedel Kentang & Rp25.397 & 397 & Daun Ubi Rebus & Rp25.397 \\ \hline
398 & Peyek Udang & Rp25.397 & 399 & Kerupuk Jangek & Rp25.397 \\ \hline
400 & Kerupuk Udang & Rp25.397 & 401 & Kerupuk Emping & Rp25.397 \\ \hline
402 & Sayur & Rp28.571 & 403 & Cabe Merah & Rp30.159 \\ \hline
404 & Gulai Telur & Rp30.159 & 405 & Telur & Rp31.746 \\ \hline
406 & Sayur Daun Ubi Gulai & Rp33.333 & 407 & Sayur Buncis & Rp33.800 \\ \hline
408 & Rendang Jengkol & Rp37.700 & 409 & Gulai Telur & Rp39.683 \\ \hline
410 & Hanya Lauk Saja & Rp39.683 & 411 & Ayam Sambal Ijo & Rp39.683 \\ \hline
412 & Nasi Ktk + Telur & Rp39.683 & 413 & Sayur Kailan & Rp40.300 \\ \hline
414 & Sayur Kangkung & Rp41.270 & 415 & Sayur Taucho & Rp41.270 \\ \hline
416 & Sayur Asem & Rp41.270 & 417 & Sayur Daun Ubi Tumbuk & Rp41.270 \\ \hline
418 & Sayur Brokoli & Rp41.270 & 419 & Ikan Gembung Acar & Rp42.857 \\ \hline
420 & Dendeng Sambal Ijo & Rp46.032 & 421 & Sambal Kerang & Rp66.667 \\ \hline
422 & Sambal Hati & Rp73.016 & 423 & Sambal Udang & Rp80.952 \\ \hline


\rowcolor{black!8}\multicolumn{6}{|l|}{\textbf{C. Minuman}} \\ \hline
424 & Teh Manis Panas & Rp25.397 & 425 & Teh Susu Panas & Rp33.333 \\ \hline
426 & Jus Tomat & Rp36.508 & 427 & Jus Wortel & Rp36.508 \\ \hline
428 & Jus Semangka & Rp36.508 & 429 & Lemon Tea & Rp36.508 \\ \hline
430 & Teh Susu Dingin & Rp36.508 & 431 & Es Doger Garuda & Rp44.444 \\ \hline
432 & Jus Kuini & Rp44.444 & 433 & Es Pelangi & Rp44.444 \\ \hline
434 & Es Alpukat Garuda & Rp44.444 & 435 & Jus Jeruk & Rp44.444 \\ \hline
436 & Jus Sunkist & Rp52.381 & 437 & Es Garuda Special & Rp52.381 \\ \hline
438 & Jus Jeruk Kelapa Muda & Rp52.381 & & & \\ \hline



\rowcolor{black!8}\multicolumn{6}{|l|}{\textbf{D. Paket Restoran Garuda X Teh Pucuk}} \\ \hline
439 & Nasi Bks Ayam Gulai + Teh Pucuk & Rp53.300 & 440 & Nasi Bks Ayam Goreng + Teh Pucuk & Rp53.300 \\ \hline
441 & Nasi Bks Rendang + Teh Pucuk & Rp57.200 & & & \\ \hline



\rowcolor{black!8}\multicolumn{6}{|l|}{\textbf{E. Paket Promo Makanan + Minuman}} \\ \hline
442 & 1 Bks Nasi Ayam Goreng + Jus Jeruk & Rp74.603 & 443 & 1 Bungkus Nasi Ayam Panggang + Jus Mangga & Rp80.600 \\ \hline
444 & 2 Bks Nasi Ikan Panggang + Jus Pokat & Rp124.800 & 445 & Kepala Kambing + Jus Timun & Rp265.851 \\ \hline
\end{xltabular}
\tablesource[\textwidth]{Observasi GrabFood kompetitor, 19--20 Juni 2026; diolah peneliti (2026).}
\endgroup

